\section{相关研究现状}

针对LLM Agent工具调用链的间接提示注入攻击威胁，现有防御工作大致沿着“提示工程与输入净化—训练与安全对齐—Agent系统架构设计”的技术路径演进，以下从这三个层面分别梳理现有研究工作的脉络及其局限。

\subsection{基于提示工程与输入净化的防御方法}

早期研究试图仅通过提示文本的组织方式或简单的净化规则抵御间接提示注入，其共同特点是既不改变模型参数，也不改变Agent的执行架构，而是依赖提示层面的结构调整引导模型区分指令与数据。Schulhoff\cite{schulhoff2024sandwich}提出的Prompt Sandwiching（提示夹心法）在不可信内容前后重复插入用户原始指令，试图借助模型对近端上下文的注意力偏好，使其在读取外部内容后仍能回忆并优先执行最初的任务目标。Hines等\cite{hines2024defending}提出的Spotlighting（数据分隔符法）通过对工具返回的不可信内容施加特殊分隔符标记或字符编码变换（如虚假语言编码、编码转写），使模型有机会将其显式识别为数据而非指令，从而降低对其中指令性内容的遵从概率。Zhang等\cite{zhang2025asb}提出的Agent Security Bench基准进一步将分隔符、提示夹心法、指令预防（Instructional Prevention）、复述（Paraphrasing）与乱序（Shuffle）等同类简单方法系统化为统一的基线家族，并纳入同一评测协议中比较其防御效果。Shi等\cite{shi2025promptarmor}提出的PromptArmor则在数据净化侧引入了更工程化的检测流程，将输入内容交由专用的外部护栏LLM做净化处理。

综合现有研究，提示工程与输入净化方法的局限主要体现在约束强度与自适应攻击鲁棒性两个方面。在约束强度方面，提示夹心法未对不可信内容进行显式标记或隔离，Spotlighting则依赖模型正确识别分隔符或编码格式，二者均以LLM对重复指令、分隔符和结构化标记的概率性遵从作为防御基础。当注入内容与任务语义高度相关，或分隔格式遭到模仿和破坏时，相应防御信号可能失效\cite{schulhoff2024sandwich,hines2024defending}。在攻击适应性方面，Agent Security Bench的评测结果显示，简单提示防御的效果会随攻击话术变化而产生明显波动\cite{zhang2025asb}；PromptArmor的语义净化依赖外部护栏LLM的识别能力与鲁棒性，仍存在检测机制被对抗性绕过的风险\cite{shi2025promptarmor}。Ji等\cite{ji2025taxonomy}对13个主流间接提示注入防御系统开展的横向自适应攻击评估进一步表明，当攻击者掌握防御机制并据此构造规避性载荷时，多数提示工程类防御的有效性较静态评测显著降低。因此，仅依靠提示层面的结构化引导尚不足以为Agent工具调用链建立确定且可验证的安全边界。

\subsection{基于训练的防御方法}

为弥补提示工程类方法的不足，已有一批研究转向通过训练或微调的方式，使大模型或专用检测器具备更稳定的指令边界识别与抗注入能力。在大模型层的安全对齐方面，Wallace等\cite{wallace2024instruction}提出的Instruction Hierarchy显式定义系统消息、用户输入与第三方工具/数据输出之间的优先级，并提出“上下文综合”与“上下文忽略”两种自动数据生成方法，训练模型在指令发生冲突时优先遵循系统级特权指令、忽略或驳回与其不一致的低特权内容。Chen等\cite{chen2025struq}提出的StruQ在LLM系统前构建前端，通过特殊标记将指令与数据分离到不同的结构化通道，并微调专门的模型识别这些标记，以降低数据被误解析为指令的概率。Chen等提出的SecAlign\cite{chen2025secalign}和Meta SecAlign\cite{chen2025metasecalign}将提示注入防御进一步形式化为偏好优化的对齐训练问题，采用直接偏好优化范式训练模型在面对含注入指令的输入时偏好安全输出、抑制不安全响应，在保持常规任务能力的同时提升对提示注入的鲁棒性。Mou等\cite{mou2026toolsafe}提出的ToolSafe框架引入基于多任务防御反馈强化学习微调的步骤级主动安全护栏模型，在每步工具调用前进行安全检查。在提示注入文本检测方面，ProtectAI\cite{protectai2024}提出基于DeBERTa-v3-base微调的提示注入分类器，以较低的推理开销对输入文本进行二分类判别；Li等\cite{li2025piguard}提出的PIGuard针对已有提示注入护栏普遍存在的“过度防御”问题进行专门优化，在维持较高攻击检出率的同时降低良性任务的误报率。Liu等\cite{liu2025datasentinel}提出的DataSentinel则采用博弈论意义下的极小极大（minimax）检测训练框架，将检测器训练建模为攻击者与检测器之间的对抗博弈，以提升检测器对适应性攻击的鲁棒性。

相较于提示工程，基于训练的防御方法能够获得较为稳定的统计安全收益，但在分布外泛化、执行约束和部署维护方面仍存在局限。从泛化能力看，Instruction Hierarchy、SecAlign与Meta SecAlign的安全性取决于参数化模型在训练分布内形成的经验性行为，尚不能保证模型在面对训练阶段未覆盖的攻击话术或自适应攻击时仍严格遵循指令层级\cite{wallace2024instruction,chen2025secalign,chen2025metasecalign}。StruQ还以分隔标记不被穿透或伪造为前提，并要求基座模型针对特定输入结构进行微调，增加了跨场景部署与迁移的成本\cite{chen2025struq}。从执行约束看，ToolSafe以提示形式向Agent反馈安全判断，反馈机制与工具执行过程之间仍为松耦合关系，难以形成对危险工具调用的强制约束\cite{mou2026toolsafe}。从检测粒度看，ProtectAI与PIGuard等外挂式检测器主要处理孤立文本片段，DataSentinel虽通过对抗训练提升了适应性攻击条件下的鲁棒性，但均未覆盖Agent动态、多轮工具调用过程中的跨步骤风险传播\cite{protectai2024,li2025piguard,liu2025datasentinel}。此外，模型训练与持续更新所需的资源投入，也制约了此类方法对快速演化攻击手段的响应能力。

\subsection{基于Agent系统架构设计的防御方法}

传统系统安全领域中，污点分析（Taint Analysis）与溯源图（Provenance Graph）长期被用于追踪不可信数据在程序执行中的传播路径、关联系统事件的因果依赖\cite{zipperle2022provenance}，为具有相似因果结构的LLM Agent工具调用链安全研究提供了方法论借鉴。然而，Agent场景下数据传播由概率性的自然语言推理而非语法依赖决定，传统方法难以直接迁移，因此近年研究开始关注运行时监控、执行溯源与因果归因以及架构级信息流控制三个方向，探索面向Agent工具调用链的专门化防御机制。

\textbf{运行时监控与执行隔离。}该类工作的防御时间节点在工具调用的实际执行环节。Shi等\cite{shi2025progent}提出的Progent权限控制框架使用运行时拦截器，在每次工具调用前通过SMT求解判定工具权限是否单调衰减，以控制权限升级类攻击。Li等\cite{li2025drift}提出的DRIFT采用基于LLM的“安全规划器生成预期工具调用轨迹—动态验证器逐步比对—注入隔离器清洗外部注入指令”三级架构，在运行时动态审批偏离预期规划的工具调用。Zhu等\cite{zhu2025melon}提出的MELON采用掩码后重执行的思路，在主Agent执行过程中引入额外的影子执行分支，将用户输入掩码替换为中性的无关指令，并通过比较掩码前后的工具调用决策判断是否受到注入攻击。Wang等\cite{wang2026agentspec}提出的AgentSpec面向运行时约束设计领域特定语言，允许开发者将安全策略转化为形式化执行规则，并在代码执行、具身Agent等高风险场景中实施拦截。Wu等\cite{wu2025isolategpt}提出的IsolateGPT将隔离机制下沉至执行环境层，为每个工具分配独立沙箱上下文，以切断工具间的状态共享与权限继承。Wen等\cite{wen2026agentsys}提出的AgentSys通过显式的分层记忆管理，在隔离的子任务执行单元内对命令类工具调用触发验证器，检测到风险后净化当前观察内容，并在原意图下有界重启。Wang等\cite{wang2025cacheprune}提出的CachePrune从模型内部表示层出发，在推理时监控并编辑KV缓存中与注入指令相关的注意力模式，引导模型不遵循被识别出的注入内容，且无需重新训练。

运行时监控与执行隔离方法在防御覆盖范围、机制适用条件和系统开销方面仍存在不同程度的限制。在防御覆盖方面，Progent主要约束权限集合的非授权扩张，对合法权限范围内实施的偏好操纵等攻击缺乏识别能力\cite{shi2025progent}；AgentSpec依赖预先编写的安全策略，AgentSys的验证器则仅覆盖隔离子任务内部的调用路径，分别受到策略扩展能力和调用入口完全中介能力的制约\cite{wang2026agentspec,wen2026agentsys}。在机制适用条件方面，DRIFT依赖安全规划器对任务轨迹的准确预判，MELON则需要通过额外的影子分支完成步骤级掩码重执行，二者均会引入附加计算过程和任务时延\cite{li2025drift,zhu2025melon}。在部署适用性方面，IsolateGPT的安全边界受沙箱隔离强度影响，CachePrune的防御机制与具体模型的缓存结构及内部表示相绑定，因而在跨模型架构或黑盒API场景中的迁移能力有限\cite{wu2025isolategpt,wang2025cacheprune}。此外，现有运行时方法多以单步“阻断或放行”的二值判定作为输出，尚难同时满足任务效用保持、细粒度响应和事后审计的需求。

\textbf{执行溯源与因果归因。}该类工作试图为Agent的运行轨迹构建可查询的溯源结构，进而定位风险的来源与传播路径。An等\cite{an2025ipiguard}提出的IPIGuard通过建模工具间的正常依赖关系，以异常依赖边识别注入攻击导致的非预期工具调用。Wang等\cite{wang2025agentarmor}提出的AgentArmor将Agent运行时轨迹抽象为程序依赖图，并引入安全类型系统进行细粒度信息流控制与策略检查。Wang等\cite{wang2026authgraph}提出的AuthGraph构建授权图与溯源图相结合的双图架构，以参数策略约束关键参数的允许来源，并将动态重规划限制在白名单与程序检查范围内。Cai等\cite{cai2026ghost}提出的NeuroTaint将关注点从语法依赖的信息流传播转向基于语义转换与因果推理的传播，通过动态上下文溯源图在Agent记忆边界和会话重启后维持污点谱系，并引入汇点驱动的因果分析器，以反事实推理捕捉隐式控制流影响。Sequeira等\cite{sequeira2026agentsentry}提出的Agent-Sentry从Agent历史良性轨迹中学习正常行为边界，构建“结构分类器初判—白名单参数来源过滤—LLM残余仲裁”的三层运行时门控架构，将大部分判定固化在确定性结构检查中。Zhang等\cite{zhang2026agentsentry}提出的AgentSentry针对工具返回边界设计受控的反事实模拟执行机制，在不提交真实副作用的前提下判断工具返回内容是否驱动后续风险行为，并通过因果门控净化上下文后继续任务。Kim等\cite{kim2026causalarmor}提出的CausalArmor对特权决策执行留一法（Leave-One-Out）意义下的归因分析，仅在观测到不可信内容对决策产生主导性影响时触发针对性净化。Yan等\cite{yan2026propguard}提出的PropGuard采用双视图时空图恢复有害传播子图，对定位到的来源节点执行抑制、净化或重新生成，并按依赖顺序重放下游交互。

执行溯源与因果归因方法能够刻画注入攻击在Agent行为轨迹中的结构指纹，但现有工作普遍存在“检测有余、闭环不足”的问题。具体而言，IPIGuard需要在执行前规划工具依赖图，对调用顺序难以预先确定的动态任务适用性不足；AgentArmor的图构建与逐调用LLM依赖边推理会增加计算开销，并可能使依赖推断模型暴露于二次注入风险\cite{an2025ipiguard,wang2025agentarmor}。AuthGraph的授权图由隔离上下文生成，尚不能独立于生成过程构成授权真值；Agent-Sentry的参数来源证据由同一Agent自行上报，来源可信度也缺乏独立验证\cite{wang2026authgraph,sequeira2026agentsentry}。NeuroTaint定位于离线事后审计，无法在执行时实时阻断或净化危险行为\cite{cai2026ghost}。AgentSentry的模拟执行不提供对象级切断回执；CausalArmor所称的“因果”仅指移除特定片段后的操作性反事实影响，且未证明净化后真实路径被永久切断；PropGuard虽提供重放验证，仍缺少切断效果的运行时物化回执、残余可达路径检查与来源非重新引入检查\cite{zhang2026agentsentry,kim2026causalarmor,yan2026propguard}。因此，此类方法尚未形成风险定位、阻断、恢复与验证的一体化闭环。

\textbf{架构级信息流控制与确定性授权。}针对运行时方法在安全保证上的概率性局限，一系列工作从系统架构层面重新设计Agent的执行模型，试图以确定性机制取代概率性过滤。Debenedetti等\cite{debenedetti2025camel}提出的CaMeL采用规划与执行分离的双LLM架构，特权规划器仅读取用户提示生成数据流图，被隔离的执行器严格按图调用工具，从架构层面阻断不可信数据对控制流的劫持。Costa等\cite{costa2025fides}提出的FIDES将信息流控制方法应用于Agent，为内容附带完整性与机密性标签并强制传播，在敏感工具执行前对不可信数据进行隔离提取与策略门控。Kim等\cite{kim2025promptflow}提出的Prompt Flow Integrity借鉴系统安全中的最小权限原则，将Agent拆分为分别处理可信与不可信数据的双组件，通过跨组件通信监控与显式、隐式流检查，在不可信数据流向敏感汇聚点时触发告警或阻断。Wu等\cite{wu2024fsecure}提出的f-secure LLM系统从信息流控制视角出发，将LLM系统解耦为具备上下文感知能力的处理管线，并依据信息流控制原则动态生成结构化执行方案，以限制不可信数据对控制流的影响。

架构级方法增强了执行约束的确定性，但代价集中在信息流覆盖能力、可信标注假设与工程可用性上。CaMeL需要额外的规划模型和隔离执行器，因而面临任务效用损失、计算成本及单次防御延时较高的问题\cite{debenedetti2025camel}。FIDES主要控制显式信息流，难以覆盖由控制决策产生的隐式流，其规划器依赖LLM进行污点数据的改写与提取，也限制了系统的开销与可扩展性\cite{costa2025fides}。Prompt Flow Integrity需要在插件层手动标注数据源并维护代理与流检查策略，存在策略规模化、任务效用和跨服务追踪方面的困难；f-secure LLM同样依赖数据来源与信任等级的准确标注，标注误差会直接传导至执行方案的安全性\cite{kim2025promptflow,wu2024fsecure}。因此，此类方法虽提供了更强的结构性约束，但其安全保证仍受到标注质量、策略完备性与部署复杂度的共同制约。

\textbf{面向授权判定与分级响应的防御机制。}在确定性信息流控制之外，另有一批工作专门关注“工具调用是否应被授权”和“判定拒绝后应如何分级响应”的闭环设计。Zhang\cite{zhang2026agentlibos}提出的Agent libOS借鉴操作系统中进程隔离与能力管理的思路，为Agent构建类似进程身份与能力授权的抽象层，将外部资源的访问权限限定在显式声明且可审计的调用接口范围内。Fan等\cite{fan2026granularity}提出的PACT从工具函数签名中自动合成参数角色契约，并在Agent跨多轮规划过程中追踪参数值的实际来源，在每次工具调用前逐参数核验来源可信度及其是否命中禁止来源规则，同时为满足特定条件的可信过程提供有限的凭证化豁免机制。Santos-Grueiro\cite{santosgrueiro2026temporary}提出的Commit-Time Authorization在工具调用即将提交执行时，将授权依据、当前派生状态与将要产生的持久效果进行绑定，核验授权的新鲜度与效果绑定的一致性后才允许提交，否则触发刷新、重新绑定、重规划或拒绝。Hossain等\cite{hossain2026nexus}提出的NEXUS在参数级检查后输出允许、阻断、确认或修订四类动作。Ma等\cite{ma2026secureclaw}提出的SecureClaw通过限定数据访问句柄与预览—提交两阶段绑定机制，在提交前由可信执行器重新核验请求的新鲜度、重放与确认状态，拒绝后仅返回粗粒度动作类别与安全提示。Sun等\cite{sun2026triad}提出的TRIAD让训练后的护栏模型在每个规划步骤输出继续、更新或拒绝三类判定并附带结构化反馈，驱动Agent在执行前修订计划。Liu\cite{liu2026toolconnectionexecutioncontrol}提出的HCP将Agent中的主体、授权、能力与数据管道等概念对象化，并检查多项执行控制不变量。

面向授权判定与分级响应的方法体现了Agent工具安全从概率性过滤走向确定性约束的趋势，但其局限可归纳为授权覆盖不足、响应校准不充分与部署验证有限。Agent libOS需要在宿主操作系统与语言运行时之上额外构建抽象层，部署与适配成本较高\cite{zhang2026agentlibos}。PACT的自动角色推断与来源追踪仍存在误差，且其控制粒度集中于单个参数，未覆盖任务级整体调用授权\cite{fan2026granularity}；Commit-Time Authorization主要保证提交状态与授权记录的一致性，并不独立判断外部内容是否构成间接提示注入\cite{santosgrueiro2026temporary}。在响应闭环方面，NEXUS在压力测试中未能正确识别应输出“修订”的40个样本，多数回退为确认动作，表明其细粒度动作校准仍有不足\cite{hossain2026nexus}；SecureClaw无法清除已经进入规划器上下文的明文内容，仍有信息层面的残余暴露\cite{ma2026secureclaw}。TRIAD仅在特定规模的训练模型、模拟工具与离线环境中得到验证，其额外时延和对未见攻击的泛化能力仍待检验；HCP的实验也局限于小规模受控场景，尚未证明能够覆盖真实部署中的多样化风险\cite{sun2026triad,liu2026toolconnectionexecutioncontrol}。总体而言，此类方法仍普遍面临部署复杂度高、任务效用折损和跨框架集成困难等挑战，且多数工作以阻断攻击为核心，缺少对风险传播路径、归因过程与残余风险的结构化解释。

综合上述工作可以看出，现有面向Agent工具调用链的提示注入攻击防御沿着提示工程、训练对齐与系统架构三条路径不断发展，呈现出两端分布的特征：提示工程与训练类方法部署成本较低，但其安全性建立在模型自身概率性遵从能力之上，在未见攻击上防御效果显著下降，难以提供鲁棒的安全性；而架构级信息流控制、可编程权限引擎与执行隔离类方法虽能提供更强的确定性约束，却普遍伴随较高的运行时开销、复杂的策略编写与维护负担，部分方法的授权判定还依赖对来源真实性或结构化账本正确性的假设，安全性与任务效用之间的矛盾尚未得到很好的调和。另一方面，现有的执行溯源、因果归因与包含性评测工作大多分别聚焦于风险检测、归因定位或状态恢复中的某一环节，尚未在同一评测框架下系统评估授权判定、跨步骤攻击链遏制与执行恢复三者的联合效果，也尚未充分探讨在调用被拒绝后向规划器披露何种粒度的反馈信息，才能在保持安全性的同时尽量维持任务效用这一问题。上述两方面的不足，为本研究分别从任务契约驱动的工具调用前防御与运行轨迹图驱动的在线检测与恢复两个研究点展开工作提供了动机。
